\section{Evaluation}

To demonstrate the effectiveness, design contribution, robustness, and ecosystem-scale practicality of \toolname, we conduct extensive experiments addressing four research questions.

\begin{itemize}[leftmargin=15pt]
    \item \textbf{RQ1: Effectiveness.}  
    How effective is \toolname when compared with state-of-the-art baselines?

    \item \textbf{RQ2: Ablation Study.}  
    How much do the major components of \toolname contribute to the overall detection performance?

    \item \textbf{RQ3: Impact of LLMs.}  
    How sensitive is \toolname to the choice of different LLMs?

    \item \textbf{RQ4: Practicality.}  
    Is \toolname practical for large-scale registries in terms of efficiency, cost, false positives, and its ability to uncover previously unreported malicious skills in the wild?
\end{itemize}

\subsection{Evaluation Setup}
\label{sec:eval-setup}

\noindent \textbf{Implementation.}
We implemented a prototype of \toolname in Python with over 7.9K lines of code, excluding third-party dependencies and external frameworks. In the parsing-based analysis, \toolname is built on top of Semgrep~\cite{semgrep}. To support realistic skill ecosystems, we incorporate a total of 2,665 Semgrep rules spanning 10 programming languages, including C, C++, C\#, Go, Java, JavaScript, PHP, Python, Ruby, and TypeScript. 
In the operand resolution, \toolname leverages the existing multi-language static analyzer YASA to perform point-to and value-flow analysis. At present, YASA supports Python, JavaScript, Java, and Go, which allows \toolname to resolve value flow and recover concrete operation objects for these languages. 
Unless otherwise specified, \toolname uses \texttt{gpt-5.3-codex-medium} as the default model in the experiments.
Due to space limitations, we do not include all prompts used in the LLM-assisted analysis in the main text. All these prompts are provided in the open-sourced artifacts.

\noindent \textbf{Running Environment.}
All experiments were conducted on a server running Ubuntu Linux 22.04.5 LTS, equipped with 256 CPUs, 1.0 TB RAM, 1 NVIDIA GeForce RTX 4090 GPU, and 28 TB of SSD storage. Unless otherwise stated, all tools were evaluated under the same environment and resource configuration.

\noindent \textbf{Dataset.}
To comprehensively evaluate \toolname, we construct two datasets, one ground truth benchmark for baseline comparison and one for large-scale in-the-wild analysis.

\textit{(1) MalSkillsBench.} 
For MalSkillsBench, we construct a labeled dataset of malicious and benign agent skills collected from real-world skill marketplaces and publicly reported malicious-skill incidents. We initially considered incorporating the malicious skills reported by MASB~\cite{liu2026maliciousagentskillswild}, which originally identified 157 malicious samples. However, a substantial portion of the corresponding GitHub repositories had become unavailable, and we were only able to successfully retrieve 57 of them. We then reran the publicly available MASB scanning pipeline on these 57 samples and found that only 6 could still be successfully detected. Moreover, these repositories are often large and artifact-heavy~(more than 4K avg.), making manual auditing and validation impractical. Since MASB is also included as one of our evaluation baselines, incorporating this small set of MASB-detected samples would additionally compromise fairness. We therefore do not include the MASB samples in the benchmark.
Instead, we construct the malicious set by systematically collecting publicly reported malicious-skill incidents from ClawHub~\cite{koi2026clawhavoc,snyk2026skills,philskents-clawhub-blocklist,opensourcemalware-clawdbot-crypto,piedpiper-openclaw-malicious-skills}. Although many of these malicious skills had already been removed from ClawHub, we recovered 597 samples from the historical commits of ClawHub's GitHub archive repository~\cite{openclaw-skills-archive}. Since many of these samples were repeatedly uploaded by the same malicious actors as part of large-scale poisoning campaigns~\cite{koi2026clawhavoc,opensourcemalware-clawdbot-crypto}, we further perform hash-based deduplication and obtain 100 unique malicious skills. Following the common practice of prior security benchmark construction~\cite{zahan2024malwarebench,cccs-cic-andmal2020,hyrum2018ember}, we build a balanced dataset by additionally selecting an equal number of popular benign skills from ClawHub. This yields a final benchmark containing 200 skills in total. For each skill, we retain its original artifact structure to preserve realistic cross-artifact analysis conditions, including prompts, manifests, scripts, configuration files, and other auxiliary artifacts.

\textit{(2) Wild-Skills-150K.}
To evaluate ecosystem-scale practicality, we further collect a large unlabeled corpus from 7 public registries and marketplaces, including 34,092 skills from Skills Directory~\cite{skillsdirectory}, 30,734 from SkillsLLM~\cite{skillsllm}, 10,598 from Smithery~\cite{smithery}, 30,789 from Skills.sh~\cite{skills.sh}, 9,643 from Skills.rest~\cite{skillsrest2026}, 11,064 from SkillsMP~\cite{skillsmp}, and 37,858 from ClawHub~\cite{clawhub,openclaw-skills-archive}. In total, this yields 164,778 raw skill entries. After deduplication, we obtain 150,108 unique skills, which we denote as \textit{Wild-Skills-150K}. This dataset is used for large-scale scanning in RQ4 to assess whether \toolname can operate efficiently and detect previously unreported malicious skills.

\noindent \textbf{Baselines.}
We compare \toolname against representative malicious skills scanning tools drawn from both industry and recent research, covering static analysis, LLM analysis, and dynamic analysis. Specifically, we include \textsc{Skill Scanner}~\cite{cisco-skill-scanner}, \textsc{Nova-Proximity}~\cite{roccia2026novaproximity}, \textsc{Skill-Sec-Scan}~\cite{huifer-skill-security-scan}, \textsc{Caterpillar}~\cite{alice2026caterpillar}, and \textsc{MASB}~\cite{liu2026maliciousagentskillswild}. To the best of our knowledge, these constitute the most representative and widely adopted publicly available end-to-end tools for malicious skill detection in the current community. We do not include Snyk's \textsc{AgentScan} in our evaluation because its detection relies on proprietary Snyk APIs and signature matching against the internal database. To ensure fair comparison, we run each baseline according to its publicly documented workflow and use its default or recommended detection settings whenever possible. For baselines that require LLM support, we use the same underlying model, \texttt{gpt-5.3-codex-medium}, across all tools. 

\subsection{RQ1: Effectiveness}
\label{sec:rq1}

We evaluate \toolname and all baselines on MalSkillsBench. For each tool, we obtain a result indicating whether a skill is malicious or benign.\footnote{For tools that output auxiliary labels such as \textit{Suspicious}, we do not count such cases as correct predictions for either the malicious or benign.} Following standard practice, we report Precision, Recall, F1-score, and False Positive Rate~(FPR).

\begin{table}[t]
\centering
\caption{Detection effectiveness on MalSkillsBench. The best results for each metric are highlighted with \textcolor{red}{red}, and the second-best results are highlighted with \textcolor{blue}{blue}. For FPR, lower values indicate better performance.}
\label{tab:rq1-main}
\resizebox{\linewidth}{!}{%
\begin{tabular}{lccrrr}
\toprule
\textbf{Tool} & \textbf{TP/FP/TN/FN} & \textbf{Prec.} & \textbf{Rec.} & \textbf{FPR} & \textbf{F1} \\
\midrule
\textsc{Skill Scanner} & 93/22/72/4 & \cellcolor{second}\textbf{0.81} & \cellcolor{best}\textbf{0.96} & 0.23 & \cellcolor{second}\textbf{0.88} \\
\textsc{Caterpillar}~\cite{alice2026caterpillar} & 82/40/60/18 & 0.67 & 0.82 & 0.40 & 0.74 \\
\textsc{MASB}~\cite{liu2026maliciousagentskillswild} & 9/10/90/91 & 0.47 & 0.09 & \cellcolor{second}\textbf{0.10} & 0.15 \\
\textsc{Skill-Sec-Scan}~\cite{huifer-skill-security-scan} & 40/50/50/60 & 0.44 & 0.40 & 0.50 & 0.42 \\
\textsc{Nova-Proximity}~\cite{roccia2026novaproximity} & 4/25/75/96 & 0.14 & 0.04 & 0.25 & 0.06 \\
\hline
\toolname & 92/5/95/8 & \cellcolor{best}\textbf{0.95} & \cellcolor{second}\textbf{0.92} & \cellcolor{best}\textbf{0.05} & \cellcolor{best}\textbf{0.93} \\
\bottomrule
\end{tabular}%
}
\end{table}

\noindent \textbf{Overall Results.}
~\autoref{tab:rq1-main} presents the overall comparison results. \toolname achieves the best overall trade-off among all compared baselines, attaining the highest F1-score of 0.93 and the lowest FPR of 0.05. Its Recall of 0.92 is slightly lower than the best-performing Recall of 0.96 achieved by \textsc{Skill Scanner}, but this small gap comes with substantially better precision and false-positive control. Notably, the overall second-best baseline is \textsc{Skill Scanner}, an industry-standard tool from CISCO that also relies on LLM-based analysis, yet it only achieves an F1-score of 0.88 with a much higher FPR of 0.23. Besides, our benchmark also shows that several tools claiming malicious-skill detection capability still perform unsatisfactorily in practice, with limited recall~(e.g., \textsc{nova-proximity} and \textsc{MASB}) or high FPR~(e.g., \textsc{Skill-Sec-Scan}). Overall, \toolname provides a better balance between precision and recall.

\noindent \textbf{FP/FN Analysis.}
We conduct a further review of the FPs and FNs. The FPs are concentrated in skills whose descriptions contain strong operational or automation-related signals, such as installation instructions, environment configuration steps, cron/heartbeat workflows, and web scraping examples. Although such content is benign in context, it partially overlaps with the capability patterns that our method treats as security-relevant evidence. 
On the other side, the FNs mainly arise in cases where maliciousness is conveyed more implicitly than our current evidence categories capture. Representative examples include documentation-centric social engineering, guidance that redirects users to external downloads, and bot-control capabilities. In these cases, the malicious signal is not always expressed as an explicit sensitive operation, but instead emerges through intent, indirection, or downstream enablement. This suggests that the main limitation of \toolname lies in covering more implicit and document-oriented threat manifestations.

\subsection{RQ2: Ablation Study}
\label{sec:rq2}

We conduct an ablation study on MalSkillsBench using the same evaluation setting as in RQ1. To understand the contribution of each major component, we evaluate the following variants:

\begin{itemize}[leftmargin=15pt]
    \item \textit{Full \toolname.} The complete system with both symbolic and neural extraction, followed by symbolic and neural reasoning.

    \item \textit{w/o Symbolic Extractor.} This variant removes the symbolic extraction component and relies only on the neural extractor to identify SSO from the skill.

    \item \textit{w/o Neuro Extractor.} This variant removes the neural extraction component and uses only the symbolic extractor for SSO collection.

    \item \textit{w/o Symbolic Reasoning.} This variant removes the symbolic reasoning stage and makes decisions based only on the remaining pipeline and neural reasoning.

    \item \textit{w/o Neuro Reasoning.} This variant removes the neural reasoning component and relies only on symbolic reasoning for final decision making.
\end{itemize}

\begin{table}[t]
\centering
\caption{Ablation study of \toolname with 4 variants.}
\label{tab:rq2-ablation}
\resizebox{\linewidth}{!}{%
\begin{tabular}{lccccc}
\toprule
\textbf{Variant} & \textbf{TP/FP/TN/FN} & \textbf{Prec.} & \textbf{Rec.} & \textbf{FPR} & \textbf{F1} \\
\midrule
\textit{Full \toolname} & 92/5/95/8 & 0.95 & 0.92 & 0.05 & 0.93 \\
\hline
\textit{w/o Symbolic Extractor} & 27/0/100/73 & 1.00 & 0.27 & 0.00 & 0.43 \\
\textit{w/o Neuro Extractor} & 91/5/95/9 & 0.95 & 0.91 & 0.05 & 0.93 \\
\textit{w/o Symbolic Reasoning} & 92/5/95/8 & 0.95 & 0.92 & 0.05 & 0.93 \\
\textit{w/o Neuro Reasoning} & 82/49/51/18 & 0.63 & 0.82 & 0.49 & 0.71 \\
\bottomrule
\end{tabular}%
}
\end{table}

\autoref{tab:rq2-ablation} presents the ablation results. Overall, the full version of \toolname achieves the best performance, with an F1-score of 0.93 and an FPR of 0.05. Among all ablations, removing the neuro reasoning component causes the largest performance drop, reducing F1 to 0.71 and increasing FPR to 0.49. This indicates that neuro reasoning is critical for accurately distinguishing malicious skills from benign ones while controlling false positives.
Removing the symbolic extractor also substantially hurts recall, which drops from 0.92 to 0.27, even though precision becomes perfect and FPR falls to 0.00. This suggests that the symbolic extractor contributes essential high-coverage evidence, and without it the system becomes overly conservative, identifying only a small subset of malicious skills.
In contrast, removing the neuro extractor leads to only a minor performance decrease, while removing symbolic reasoning yields no observable degradation on this benchmark. This is consistent with the intended roles of these two components. The neuro extractor is mainly designed to capture more stealthy and adversarial SSOs, and such benefits may not be fully reflected on the current benchmark. Meanwhile, symbolic reasoning is primarily intended to leverage structured intermediate evidence to constrain the reasoning process and reduce unnecessary token consumption, improving efficiency and reasoning stability rather than directly boosting detection metrics in every setting.

\subsection{RQ3: Impact of LLMs}
\label{sec:rq3}

\begin{table}[t]
\centering
\caption{Performance of \toolname with different LLMs.}
\label{tab:rq3-llm}
\resizebox{\linewidth}{!}{%
\begin{tabular}{lccccc}
\toprule
\textbf{LLM} & \textbf{TP/FP/TN/FN} & \textbf{Prec.} & \textbf{Rec.} & \textbf{FPR} & \textbf{F1} \\
\midrule
\textit{GPT-5.3-CodeX-Medium} & 92/5/95/8 & \cellcolor{best}\textbf{0.95} & \cellcolor{second}\textbf{0.92} & \cellcolor{best}\textbf{0.05} & \cellcolor{best}\textbf{0.93} \\
\hline
\textit{Claude-Sonnet-4.6} & 70/14/86/30 & 0.83 & 0.70 & 0.14 & 0.76 \\
\textit{Gemini-3.1-Flash-Lite-Preview} & 95/24/76/5 & 0.80 & \cellcolor{best}\textbf{0.95} & 0.24 & \cellcolor{second}\textbf{0.87} \\
\textit{Qwen3.5-397B-A17B} & 78/15/85/22 & 0.84 & 0.78 & 0.15 & 0.81 \\
\textit{DeepSeek-V3.2} & 71/5/95/29 & \cellcolor{second}\textbf{0.93} & 0.71 & \cellcolor{best}\textbf{0.05} & 0.81 \\
\bottomrule
\end{tabular}%
}
\end{table}

To understand how the choice of LLM affects \toolname, we instantiate the framework with several widely used models under the same experimental setting, including Claude Sonnet 4.6, Gemini 3.1 Flash-Lite Preview, Qwen3.5-397B-A17B, and DeepSeek-V3.2. \autoref{tab:rq3-llm} reports the results.
Overall, the choice of LLM has a clear impact on the detection performance of \toolname. GPT-5.3-CodeX-Medium achieves the best overall results, with the highest precision (0.95) and F1-score (0.93), while also maintaining a low false positive rate (0.05) and high recall (0.92). Gemini achieves the highest recall (0.95), correctly identifying most malicious skills, but at the cost of the highest false positive rate (0.24). In contrast, DeepSeek attains high precision (0.93) and the lowest false positive rate (0.05), but its recall drops to 0.71, indicating a more conservative behavior that misses more malicious samples. Qwen provides a relatively balanced middle ground (precision 0.84, recall 0.78, F1 0.81), while Claude performs slightly worse overall, particularly in recall (0.70).
Nevertheless, \toolname remains effective across all tested LLMs, achieving F1-scores ranging from 0.76 to 0.93. This suggests that while stronger LLMs can noticeably improve performance, the effectiveness of \toolname does not depend on a single model family. 

\subsection{RQ4: Practicality}
